\begin{proofbox}
\textbf{方法一：选顶点算常数（最直接）。}
若已由焦准定义或第二定义知道 $PF/PH$ 在曲线上恒定，便不必对一般点重复计算。
横向椭圆或双曲线取 $A=(a,0)$、对应 $F=(c,0)$、$l:x=a^2/c$。
\begin{center}\begin{tikzpicture}[font=\small,>=stealth]
\draw[black!30,rounded corners=2pt] (0,0) rectangle (14.4,8.2);
\node[anchor=west,font=\small\bfseries] at (.5,7.72) {焦准比若已知恒定，选实轴顶点 $A$ 就能最快算出常数};
\begin{scope}[shift={(2,5.6)},scale=.95]
\draw[black!45] (0,0)--(9,0);
\draw[violet,very thick] (8.333333333333334,-.7)--(8.333333333333334,.7);
\draw[orange!85!black,very thick] (3,.09)--(5,.09);
\fill (0,0) circle (2pt) node[below=5pt] {$O$};
\fill (5,0) circle (2pt) node[above=5pt] {$A(5,0)$};
\fill (3,0) circle (2pt) node[above=5pt] {$F(3,0)$};
\node[above=4pt] at (8.333333333333334,.7) {$l:x=25/3$};
\end{scope}
\node[anchor=west] at (1,4.7) {椭圆：$AF=2$，$d(A,l)=10/3$，$AF/d(A,l)=3/5$};
\begin{scope}[shift={(2,2.9)},scale=.95]
\draw[black!45] (0,0)--(9,0);
\draw[violet,very thick] (1.8,-.7)--(1.8,.7);
\draw[orange!85!black,very thick] (3,.09)--(5,.09);
\fill (0,0) circle (2pt) node[below=5pt] {$O$};
\fill (3,0) circle (2pt) node[above=5pt] {$A(3,0)$};
\fill (5,0) circle (2pt) node[above=5pt] {$F(5,0)$};
\node[above=4pt] at (1.8,.7) {$l:x=9/5$};
\end{scope}
\node[anchor=west] at (1,2.0) {双曲线：$AF=2$，$d(A,l)=6/5$，$AF/d(A,l)=5/3$};

\node[anchor=west] at (.7,.55) {$AF=|a-c|$；$d(A,l)=|a-a^2/c|=(a/c)|a-c|$；约去非零量};
\end{tikzpicture}\end{center}

三点共线，直接读出
\[
 e=\frac{|a-c|}{|a-a^2/c|}
 =\frac{|a-c|}{(a/c)|a-c|}=\boxed{\frac ca}.
\]
这里 $a,c>0$ 且 $a\ne c$，所以约分合法。
\textbf{此法求的是已知恒定比值的数值；它不单独证明一般点的比值恒定。}
下面给出独立代数核验，避免循环引用。
\end{proofbox}
\clearpage
\begin{proofbox}
\textbf{方法二：一般点配方，两类曲线一次写清。}
横向标准椭圆有 $y^2=b^2(1-x^2/a^2)$、$c^2=a^2-b^2$；
横向标准双曲线有 $y^2=b^2(x^2/a^2-1)$、$c^2=a^2+b^2$。
两种情况下分别代入并展开，结果相同：
\[
 \begin{aligned}
 PF^2&=(x-c)^2+y^2\\
 &=\frac{c^2}{a^2}x^2-2cx+a^2
 =\left(a-\frac ca x\right)^2.
 \end{aligned}
\]
因此无须猜测左右支的符号，统一取绝对值：
\[
 PF=\left|a-\frac ca x\right|
 =\frac ca\left|x-\frac{a^2}{c}\right|
 =\frac ca\,d(P,l).
\]
分母不能为零：椭圆上 $x\le a<a^2/c$；双曲线上 $|x|\ge a>a^2/c>0$，所以均不在这条准线上。
约去准距便得 $e=c/a$，覆盖椭圆全部点和双曲线两支。
左焦点、左准线由对称性得到；纵向交换坐标。

\medskip
\textbf{范围与形状。} 椭圆 $0<c<a$，故 $0<e<1$；双曲线 $c>a$，故 $e>1$。
抛物线由定义 $PF=PH>0$ 得 $e=1$。
两种参数平方关系除以 $a^2$，就得到结论中的根式公式。
\begin{center}\begin{tikzpicture}[font=\small,>=stealth]
\draw[black!30,rounded corners=2pt] (0,0) rectangle (14.4,8.2);
\node[anchor=west,font=\small\bfseries] at (.5,7.72) {固定长半轴 $a=3$：$e$ 增大，$b/a=\sqrt{1-e^2}$ 减小};
\begin{scope}[shift={(2.4,4)},scale=.52]
\draw[blue,very thick] (0,0) ellipse (3 and 3.0000000000);
\draw[black!35] (-3.3,0)--(3.3,0);
\fill (0.0000000000,0) circle (3pt);\fill (0.0000000000,0) circle (3pt);
\end{scope}
\node at (2.4,1.4) {$e=0$};
\begin{scope}[shift={(7.2,4)},scale=.52]
\draw[blue,very thick] (0,0) ellipse (3 and 2.4000000000);
\draw[black!35] (-3.3,0)--(3.3,0);
\fill (1.8000000000,0) circle (3pt);\fill (-1.8000000000,0) circle (3pt);
\end{scope}
\node at (7.2,1.4) {$e=0.6$};
\begin{scope}[shift={(12,4)},scale=.52]
\draw[blue,very thick] (0,0) ellipse (3 and 0.9367496998);
\draw[black!35] (-3.3,0)--(3.3,0);
\fill (2.8500000000,0) circle (3pt);\fill (-2.8500000000,0) circle (3pt);
\end{scope}
\node at (12,1.4) {$e=0.95$};

\node[anchor=west] at (.7,.55) {同尺度；$e=0$ 是圆的边界，焦点合到中心，不能使用有限准线公式 $a/e$};
\end{tikzpicture}\end{center}

固定长半轴时，$e$ 越大，$b$ 越小，椭圆越扁。
$e=0$ 是圆的边界；$e\to1^-$ 时 $b\to0$ 是椭圆扁向线段的极限，\textbf{不能把固定 $a$ 的这个极限直接说成抛物线}。
双曲线则有 $b/a=\sqrt{e^2-1}$，其朝实轴正方向的渐近线夹角随 $e$ 增大。
\end{proofbox}
\clearpage
